% ============================================================================
%  OhMyGAD — Admin User Manual
%  LaTeX Template
% ============================================================================
\documentclass[12pt, a4paper]{article}

% ── Packages ────────────────────────────────────────────────────────────────
\usepackage[utf8]{inputenc}
\usepackage[T1]{fontenc}
\usepackage{lmodern}
\usepackage[margin=1in]{geometry}
\usepackage{graphicx}
\usepackage{xcolor}
\usepackage{hyperref}
\usepackage{booktabs}
\usepackage{longtable}
\usepackage{enumitem}
\usepackage{fancyhdr}
\usepackage{titlesec}
\usepackage{tocloft}
\usepackage{float}
\usepackage{caption}
\usepackage{parskip}
\usepackage{tabularx}
\usepackage{array}

% ── Colours ─────────────────────────────────────────────────────────────────
\definecolor{primarydark}{HTML}{2D2A4A}
\definecolor{periwinkle}{HTML}{B8B5E8}
\definecolor{softpink}{HTML}{F4A7B9}
\definecolor{accentgold}{HTML}{F4C97A}
\definecolor{linkblue}{HTML}{2563EB}
\definecolor{lightgray}{HTML}{F5F5F5}

% ── Hyperlink styling ──────────────────────────────────────────────────────
\hypersetup{
  colorlinks  = true,
  linkcolor   = primarydark,
  urlcolor    = linkblue,
  citecolor   = primarydark,
  pdfauthor   = {OhMyGAD Development Team},
  pdftitle    = {OhMyGAD Admin User Manual},
  pdfsubject  = {User Manual for Admin Users},
}

% ── Header / Footer ───────────────────────────────────────────────────────
\pagestyle{fancy}
\fancyhf{}
\fancyhead[L]{\small\textcolor{primarydark}{OhMyGAD — Admin User Manual}}
\fancyhead[R]{\small\textcolor{primarydark}{\thepage}}
\fancyfoot[C]{\small\textcolor{gray}{Confidential — For Internal Use Only}}
\renewcommand{\headrulewidth}{0.4pt}
\renewcommand{\footrulewidth}{0.2pt}

% ── Section styling ────────────────────────────────────────────────────────
\titleformat{\section}
  {\Large\bfseries\color{primarydark}}{\thesection}{1em}{}
\titleformat{\subsection}
  {\large\bfseries\color{primarydark}}{\thesubsection}{1em}{}
\titleformat{\subsubsection}
  {\normalsize\bfseries\color{primarydark}}{\thesubsubsection}{1em}{}

% ── Reusable screenshot placeholder ───────────────────────────────────────
\newcommand{\screenshotplaceholder}[1]{%
  \begin{figure}[H]
    \centering
    \fbox{\parbox{0.85\textwidth}{\vspace{4cm}
      \centering\textcolor{gray}{\textit{Screenshot: #1}}
    \vspace{4cm}}}
    \caption{#1}
  \end{figure}
}

% ── Reusable field-description table ──────────────────────────────────────
% Usage: \fieldtable{ rows with \fieldrow }
\newenvironment{fieldtable}{%
  \begin{longtable}{|>{\raggedright\arraybackslash}p{2.8cm}|>{\raggedright\arraybackslash}p{1.4cm}|>{\raggedright\arraybackslash}p{2.2cm}|>{\raggedright\arraybackslash}p{6cm}|}
  \hline
  \rowcolor{lightgray}
  \textbf{Field} & \textbf{Required} & \textbf{Limit} & \textbf{Validation \& Notes} \\
  \hline
  \endfirsthead
  \hline
  \rowcolor{lightgray}
  \textbf{Field} & \textbf{Required} & \textbf{Limit} & \textbf{Validation \& Notes} \\
  \hline
  \endhead
  \hline
  \endfoot
}{%
  \end{longtable}
}
\newcommand{\fieldrow}[4]{%
  #1 & #2 & #3 & #4 \\ \hline
}

% ════════════════════════════════════════════════════════════════════════════
\begin{document}

% ── Title Page ─────────────────────────────────────────────────────────────
\begin{titlepage}
  \centering
  \vspace*{3cm}

  % TODO: Replace with actual project logo
  % \includegraphics[width=0.3\textwidth]{logo.png}\\[1.5cm]

  {\Huge\bfseries\textcolor{primarydark}{OhMyGAD}}\\[0.4cm]
  {\LARGE\textcolor{primarydark}{Admin User Manual}}\\[1.5cm]

  {\large Gender and Development Information System}\\[0.5cm]
  {\large University of the Philippines Baguio}\\[2cm]

  \textcolor{gray}{\today}\\[0.5cm]
  \textcolor{gray}{Version 1.0}

  \vfill
  {\small Prepared by the OhMyGAD Development Team}
\end{titlepage}

% ── Table of Contents ──────────────────────────────────────────────────────
\tableofcontents
\newpage

% ════════════════════════════════════════════════════════════════════════════
%  1. INTRODUCTION
% ════════════════════════════════════════════════════════════════════════════
\section{Introduction}

\subsection{What is OhMyGAD?}
OhMyGAD is a web-based information system built for the \textbf{Gender and Development (GAD)} office of UPB. It helps admin staff manage events, surveys, guidelines, and user accounts — all from one place.

\subsection{Who Is This Manual For?}
This manual is written for users who have been given the \textbf{Admin} role. Admins have the highest level of access and can manage all parts of the system, including creating and editing events, surveys, guidelines, and user accounts.

\subsection{How to Use This Manual}
\begin{itemize}[leftmargin=1.5em]
  \item Each section covers one major page or feature of the system.
  \item \textbf{Screenshots} are provided to help you identify what you see on screen.
  \item \textbf{Field tables} list every input you can fill in, along with its character limit and what the system checks for (validation rules).
  \item Fields marked with an asterisk (\textbf{*}) are required — the form will not submit without them.
\end{itemize}

\subsection{Accessing the System}
\begin{enumerate}[leftmargin=1.5em]
  \item Open your web browser (Chrome, Firefox, or Edge recommended).
  \item Navigate to the system URL provided by your administrator.
  \item Log in using your email and password, or sign in with Google.
\end{enumerate}

\screenshotplaceholder{Login Page}

% ════════════════════════════════════════════════════════════════════════════
%  2. DASHBOARD
% ════════════════════════════════════════════════════════════════════════════
\newpage
\section{Dashboard}

After logging in you will land on the \textbf{Dashboard}. This page gives you a quick summary of what is happening across the system.

\textbf{What you will see:}
\begin{itemize}[leftmargin=1.5em]
  \item Quick statistics (number of users, upcoming events, open surveys, etc.).
  \item Shortcuts to common actions.
\end{itemize}

\screenshotplaceholder{Admin Dashboard — Overview}

\subsection{Global Search}
At the top of every page you will find the \textbf{Search bar}. You can use it to quickly find events, users, guidelines, or surveys by typing a keyword.

\begin{fieldtable}
  \fieldrow{Search query}{No}{64 characters}{Type at least 1 character to see results. Press Enter to open the full search page.}
\end{fieldtable}

\screenshotplaceholder{Global Search — Dropdown Results}

% ════════════════════════════════════════════════════════════════════════════
%  3. EVENTS MANAGEMENT
% ════════════════════════════════════════════════════════════════════════════
\newpage
\section{Events Management}

The \textbf{Events} page lets you view, create, edit, and track all GAD-related events such as orientations, forums, research sessions, training, and workshops.

\screenshotplaceholder{Events Management — List View}

% ── 3.1 Creating / Editing an Event ───────────────────────────────────────
\subsection{Creating or Editing an Event}

Click the \textbf{``Create Event''} button (or the edit icon on an existing event) to open the event form. The form is split into two sections: \textit{Basic Information} and \textit{Event Schedule}.

\screenshotplaceholder{Event Form — Basic Information (Left Column)}
\screenshotplaceholder{Event Form — Event Schedule (Right Column)}

\subsubsection{Basic Information Fields}

\begin{fieldtable}
  \fieldrow{Title *}{Yes}{500 chars}{Cannot be empty. Briefly describe the event (e.g.\ ``Gender Sensitivity Orientation'').}
  \fieldrow{Description}{No}{5\,000 chars}{If provided, must be at least 10 characters long. Describe the purpose and agenda of the event.}
  \fieldrow{Location *}{Yes}{100 chars}{Cannot be empty. Enter the venue name (e.g.\ ``Sarmiento Hall'').}
  \fieldrow{Banner Image}{No}{5\,MB file}{Accepted formats: JPG, PNG, or WebP. This image appears at the top of the event detail view.}
  \fieldrow{Capacity *}{Yes}{1 -- 500}{Must be a whole number greater than 0 and at most 500. Negative numbers and decimals are not allowed.}
  \fieldrow{Category *}{Yes}{Dropdown}{Choose one: GSO, ASHO, Forum, Research, Training, or Workshop. The placeholder ``Select category'' cannot be re-selected once a choice is made.}
  \fieldrow{Status}{Read-only}{---}{Automatically calculated from the start and end dates. Possible values: \textit{Upcoming}, \textit{Today}, or \textit{Past}.}
\end{fieldtable}

\subsubsection{Event Schedule Fields}

\begin{fieldtable}
  \fieldrow{Start Date \& Time *}{Yes}{Date-time picker}{Must be set in the future (for new events). The event cannot start after it ends.}
  \fieldrow{End Date \& Time *}{Yes}{Date-time picker}{Must be after the Start Date. The minimum selectable date is automatically set to the chosen start date.}
  \fieldrow{Registration Opens *}{Yes}{Date-time picker}{Must be before the event ends. This is when users can begin registering.}
  \fieldrow{Registration Closes *}{Yes}{Date-time picker}{Must be after Registration Opens, must be before or at the event end date, and must be before the event starts. The minimum selectable date is the Registration Opens date.}
\end{fieldtable}

\subsubsection{Date Validation Summary}
The system enforces several rules to make sure your dates make sense:
\begin{enumerate}[leftmargin=1.5em]
  \item End date must be \textbf{after} the start date.
  \item Registration close must be \textbf{after} registration open.
  \item Registration must open \textbf{before} the event ends.
  \item Registration must close \textbf{before or when} the event ends.
\end{enumerate}

If any rule is broken, a red error message will appear below the relevant date field.

\screenshotplaceholder{Event Form — Date Validation Error Example}

% ════════════════════════════════════════════════════════════════════════════
%  4. SURVEYS MANAGEMENT
% ════════════════════════════════════════════════════════════════════════════
\newpage
\section{Surveys Management}

The \textbf{Surveys} page lets you create post-event feedback forms, view responses, and analyze results.

\screenshotplaceholder{Surveys Management — List View}

\subsection{Creating or Editing a Survey}

Click \textbf{``Create Survey''} to open the survey form. It has three sections: \textit{Survey Details}, \textit{Availability}, and \textit{Questions}.

\screenshotplaceholder{Survey Form — Details and Availability}

\subsubsection{Survey Details Fields}

\begin{fieldtable}
  \fieldrow{Title *}{Yes}{500 chars}{Cannot be empty (e.g.\ ``Post-Event Feedback Form'').}
  \fieldrow{Description}{No}{3\,000 chars}{If provided, must be at least 5 characters long.}
  \fieldrow{Linked Event *}{Yes}{Dropdown}{Select the event this survey is for. The placeholder ``Select linked event'' is disabled once a choice is made. Long event names are automatically shortened in the dropdown.}
  \fieldrow{Status}{Read-only}{---}{Automatically calculated: \textit{Upcoming} (before open), \textit{Open} (currently accepting responses), or \textit{Closed} (past close date).}
\end{fieldtable}

\subsubsection{Availability Fields}

\begin{fieldtable}
  \fieldrow{Opens At *}{Yes}{Date-time picker}{When the survey becomes available to respondents. Must be in the future for new surveys. Minimum selectable date is the linked event's end date.}
  \fieldrow{Closes At *}{Yes}{Date-time picker}{When the survey stops accepting responses. Must be after the Opens At date.}
\end{fieldtable}

\subsubsection{Questions Section}

You can add up to \textbf{50 questions} per survey. Each question has the following fields:

\begin{fieldtable}
  \fieldrow{Question Text *}{Yes}{500 chars}{The question you want to ask (e.g.\ ``How would you rate this event?'').}
  \fieldrow{Question Type}{Yes}{Dropdown}{Choose one: \textit{Text (open-ended)}, \textit{Multiple Choice}, \textit{Likert Scale (1--5)}, or \textit{Yes / No}.}
  \fieldrow{Choices}{Conditional}{Up to 10 choices, 100 chars each}{Only shown for Multiple Choice questions. At least 2 non-empty choices are required. Duplicate choices are not allowed.}
  \fieldrow{Required toggle}{No}{Checkbox}{If checked, respondents must answer this question before submitting.}
\end{fieldtable}

\screenshotplaceholder{Survey Form — Questions Section}

\subsection{Survey Analytics}

After a survey receives responses, you can view analytics by clicking on the survey. The analytics modal shows each question's responses, including charts for multiple-choice and rating questions.

\screenshotplaceholder{Survey Analytics Modal}

% ════════════════════════════════════════════════════════════════════════════
%  5. GUIDELINES MANAGEMENT
% ════════════════════════════════════════════════════════════════════════════
\newpage
\section{Guidelines Management}

The \textbf{Guidelines} page lets you create and manage informational documents (e.g.\ ``Gender and Technology'') that are visible to all users.

\screenshotplaceholder{Guidelines Management — List View}

\subsection{Creating or Editing a Guideline}

\screenshotplaceholder{Guideline Form}

\begin{fieldtable}
  \fieldrow{Title *}{Yes}{500 chars}{Cannot be empty (e.g.\ ``Gender and Technology'').}
  \fieldrow{Description}{No}{100\,000 chars}{A long-form text area for the guideline content. This supports detailed, multi-paragraph content.}
\end{fieldtable}

% ════════════════════════════════════════════════════════════════════════════
%  6. USER MANAGEMENT
% ════════════════════════════════════════════════════════════════════════════
\newpage
\section{User Management}

The \textbf{Users} page lets you view all registered users, create new accounts, and edit existing ones. This page is \textbf{only available to Admin users}.

\screenshotplaceholder{User Management — List View}

\subsection{Creating or Editing a User}

The user form can appear in two layouts: a \textbf{modal} (pop-up) for quick edits, or a \textbf{full page} for detailed entry. Both contain the same fields.

\screenshotplaceholder{User Form — Account Information Section}

\subsubsection{Account Information Fields}

\begin{fieldtable}
  \fieldrow{Full Name *}{Yes}{64--100 chars}{Cannot be empty. Only letters, spaces, periods, apostrophes, hyphens, and underscores are allowed.}
  \fieldrow{Email *}{Yes}{50 chars}{Must be a valid email ending in \texttt{@gmail.com} or \texttt{@up.edu.ph}.}
  \fieldrow{Password *}{Yes (create); No (edit)}{8--128 chars}{Required when creating a new user. When editing, leave blank to keep the current password. Must be at least 8 characters.}
  \fieldrow{Role *}{Yes}{Dropdown}{Choose one: Admin, Staff, Faculty, or Student. Changing the role on an existing user will trigger a confirmation prompt and reset role-specific fields.}
\end{fieldtable}

\subsubsection{Personal Information Fields}

\begin{fieldtable}
  \fieldrow{Display Name}{No}{30 chars}{A nickname or preferred name. Only letters, spaces, and basic punctuation are allowed.}
  \fieldrow{Contact Number}{No}{Exactly 11 digits}{Only digits are allowed (e.g.\ ``09123456789'').}
  \fieldrow{Address}{No}{64--100 chars}{Free-text field for the user's address.}
\end{fieldtable}

\subsubsection{Role-Specific Fields}

Depending on the selected role, additional fields will appear:

\paragraph{Student-Only Fields}

\begin{fieldtable}
  \fieldrow{Student Number *}{Yes}{Exactly 9 digits}{Only digits allowed. The system automatically derives the year level from the first 4 digits (admission year).}
  \fieldrow{Year Level *}{Yes}{Dropdown}{Options: 1st--5th Year, or Extendee. Automatically set when a valid student number is entered. A warning appears if the selected year does not match the derived year.}
  \fieldrow{College *}{Yes}{Dropdown}{Choose one: College of Science (CS), College of Arts and Communication (CAC), or College of Social Sciences (CSS).}
  \fieldrow{Program *}{Yes}{Dropdown}{Options depend on the selected college. You must select a college first.}
\end{fieldtable}

\paragraph{Faculty-Only Fields}

\begin{fieldtable}
  \fieldrow{College}{No}{Dropdown}{Same options as students.}
  \fieldrow{Department}{No}{64 chars}{Free-text (e.g.\ ``Dept.\ of Math and Computer Science'').}
\end{fieldtable}

\paragraph{Admin / Staff Fields}

\begin{fieldtable}
  \fieldrow{Office / Unit}{No}{64 chars}{Free-text (e.g.\ ``Office of the Chancellor'').}
\end{fieldtable}

\subsubsection{Identity Fields (All Roles)}

\screenshotplaceholder{User Form — Identity and Session Attendance}

\begin{fieldtable}
  \fieldrow{Pronouns}{No}{Dropdown}{Options: he/him, she/her, they/them, he/they, she/they, any/all, or Prefer not to say.}
  \fieldrow{Sex at Birth *}{Yes}{Dropdown}{Options: Male, Female, Intersex, or Prefer not to say.}
  \fieldrow{Gender Identity *}{Yes}{Dropdown}{Options: Man, Woman, Non-binary, Genderqueer, Genderfluid, or Prefer not to say.}
\end{fieldtable}

\subsubsection{Session Attendance Fields (Students \& Faculty)}

These fields track how many orientation/session events a user has attended. Each field accepts a whole number from \textbf{0 to 5}.

\begin{fieldtable}
  \fieldrow{GSO Sessions}{No}{0--5}{Gender Sensitivity Orientation sessions attended.}
  \fieldrow{ASHO Sessions}{No}{0--5}{Anti-Sexual Harassment Orientation sessions attended.}
  \fieldrow{Forums}{No}{0--5}{Forum sessions attended.}
  \fieldrow{Research}{No}{0--5}{Research sessions attended.}
  \fieldrow{Training}{No}{0--5}{Training sessions attended.}
  \fieldrow{Workshops}{No}{0--5}{Workshop sessions attended.}
\end{fieldtable}

\subsubsection{Onboarding Toggle}

\begin{fieldtable}
  \fieldrow{Onboarding complete}{No}{Toggle (on/off)}{Indicates whether the user has finished the onboarding process. Defaults to \textit{on}. Turn it off to force the user through onboarding again on their next login.}
\end{fieldtable}

% ════════════════════════════════════════════════════════════════════════════
%  7. PROFILE PAGE
% ════════════════════════════════════════════════════════════════════════════
\newpage
\section{Profile Page}

Your \textbf{Profile} page lets you view and update your own personal information. It is organized into three tabs: \textit{Personal}, \textit{Professional}, and \textit{Identity}.

\screenshotplaceholder{Admin Profile Page — Personal Tab}

\subsection{Personal Tab}

\begin{fieldtable}
  \fieldrow{Full Name}{Yes}{64 chars}{Only letters, spaces, and basic punctuation. Cannot be empty.}
  \fieldrow{Display Name}{No}{32 chars}{Only letters, spaces, and basic punctuation.}
  \fieldrow{Pronouns}{No}{Dropdown}{Options: he/him, she/her, they/them, he/they, she/they, any/all, or Prefer not to say.}
  \fieldrow{Contact Number}{No}{11 digits}{Only digits (e.g.\ ``09171234567''). Must be exactly 11 digits if provided.}
  \fieldrow{Address}{No}{100 chars}{Free-text field.}
\end{fieldtable}

\subsection{Professional Tab}

\begin{fieldtable}
  \fieldrow{Office / Unit}{No}{64 chars}{Your office or unit within the university (e.g.\ ``Office of the Chancellor'').}
\end{fieldtable}

\screenshotplaceholder{Admin Profile Page — Professional Tab}

\subsection{Identity Tab}

\begin{fieldtable}
  \fieldrow{Sex at Birth}{No}{Dropdown}{Options: Male, Female, Intersex, or Prefer not to say.}
  \fieldrow{Gender Identity}{No}{Dropdown}{Options: Man, Woman, Non-binary, Genderqueer, Genderfluid, or Prefer not to say.}
\end{fieldtable}

\textbf{Privacy Note:} The identity information is kept strictly private and is used only for institutional reporting. You are not required to fill this out.

\screenshotplaceholder{Admin Profile Page — Identity Tab}

\subsection{Profile Sidebar}

On the left side of the profile page, you will see a summary card showing:
\begin{itemize}[leftmargin=1.5em]
  \item Your name and display name.
  \item Your role badge (``Administrator'') and office.
  \item Progress bars for session attendance (GSO, ASHO, Forums, Research, Training, Workshops).
\end{itemize}

\screenshotplaceholder{Admin Profile Page — Sidebar Summary Card}

% ════════════════════════════════════════════════════════════════════════════
%  8. SETTINGS PAGE
% ════════════════════════════════════════════════════════════════════════════
\newpage
\section{Settings Page}

The \textbf{Settings} page lets you update your login credentials. It has two sections: \textit{Change Email Address} and \textit{Change Password}.

\screenshotplaceholder{Settings Page}

\subsection{Change Email Address}

\begin{fieldtable}
  \fieldrow{Current Email}{Read-only}{---}{Your current email is shown but cannot be edited directly.}
  \fieldrow{New Email *}{Yes}{No explicit limit}{Must be a valid email address. After submitting, verification links will be sent to both your old and new email addresses. You must confirm from both inboxes.}
\end{fieldtable}

\subsection{Change Password}

\textbf{Note:} This section is only available if you signed up with an email and password. If you use Google Sign-In, a message will explain that password management is handled through your Google account.

\begin{fieldtable}
  \fieldrow{Current Password *}{Yes}{16 chars}{Your existing password. The system verifies it before allowing a change.}
  \fieldrow{New Password *}{Yes}{16 chars}{Must be at least 8 characters long.}
  \fieldrow{Confirm Password *}{Yes}{16 chars}{Must match the New Password exactly.}
\end{fieldtable}

% ════════════════════════════════════════════════════════════════════════════
%  9. TIPS AND TROUBLESHOOTING
% ════════════════════════════════════════════════════════════════════════════
\newpage
\section{Tips and Troubleshooting}

\subsection{Common Validation Errors}

\begin{longtable}{|>{\raggedright\arraybackslash}p{5cm}|>{\raggedright\arraybackslash}p{8cm}|}
\hline
\rowcolor{lightgray}
\textbf{Error Message} & \textbf{What To Do} \\
\hline
``Full name is required.'' & Make sure the Full Name field is not blank. \\
\hline
``Full name can only contain letters, spaces, and basic punctuation.'' & Remove any numbers or special characters (e.g.\ \texttt{\#}, \texttt{\&}) from the name. \\
\hline
``Email must end with @gmail.com or @up.edu.ph.'' & Only these two email domains are accepted. \\
\hline
``Password must be at least 8 characters long.'' & Use a longer password with at least 8 characters. \\
\hline
``Contact number must be 11 digits.'' & Enter exactly 11 digits with no spaces or dashes. \\
\hline
``Student number must be exactly 9 digits.'' & Enter the full 9-digit student number. \\
\hline
``End date must be after the start date.'' & Adjust your event dates so the end comes after the start. \\
\hline
``Registration close must be after open.'' & Make sure the registration closing date/time is later than the opening date/time. \\
\hline
``At least one question is required.'' & Add at least one question to your survey before submitting. \\
\hline
``Question X needs at least 2 non-empty choices.'' & Multiple-choice questions need at least 2 answer options that are not blank. \\
\hline
\end{longtable}

\subsection{General Tips}
\begin{itemize}[leftmargin=1.5em]
  \item \textbf{Save often.} Changes are only saved when you click the Save or Submit button.
  \item \textbf{Check for red error messages.} They appear directly below the field that has a problem.
  \item \textbf{Dropdown placeholders are disabled.} Once you pick an option from a dropdown (e.g.\ ``Select category''), you cannot go back to the placeholder. Pick a valid option instead.
  \item \textbf{Use the search bar.} It is the fastest way to find any event, user, guideline, or survey.
  \item \textbf{Date pickers enforce minimum dates.} For example, the End Date picker will not let you pick a date before the Start Date.
\end{itemize}

% ════════════════════════════════════════════════════════════════════════════
%  10. GLOSSARY
% ════════════════════════════════════════════════════════════════════════════
\newpage
\section{Glossary}

\begin{longtable}{|>{\raggedright\arraybackslash}p{4cm}|>{\raggedright\arraybackslash}p{9cm}|}
\hline
\rowcolor{lightgray}
\textbf{Term} & \textbf{Definition} \\
\hline
GAD & Gender and Development — the university office focused on gender equity programs. \\
\hline
GSO & Gender Sensitivity Orientation — a mandatory orientation session. \\
\hline
ASHO & Anti-Sexual Harassment Orientation — a mandatory orientation session. \\
\hline
Onboarding & The initial setup process a new user completes after their first login (selecting college, program, identity info, etc.). \\
\hline
Likert Scale & A 1-to-5 rating scale commonly used in surveys (1 = Strongly Disagree, 5 = Strongly Agree). \\
\hline
Banner Image & The cover photo shown at the top of an event's detail view. \\
\hline
\end{longtable}

\end{document}
